\documentclass[11pt,letterpaper]{article}

\usepackage[margin=0.6in]{geometry}
\usepackage{enumitem}
\usepackage{amsmath}
\usepackage{tikz}
\usetikzlibrary{decorations.pathmorphing, decorations.pathreplacing, calc, arrows.meta}

\pagestyle{empty}
\setlength{\parindent}{0pt}

\begin{document}

\begin{center}
    {\Large \textbf{16.07 Dynamics --- Quiz 2: Foucault's Pendulum}}
\end{center}

\vspace{4pt}

Name: \rule{3.2in}{0.4pt}

\vspace{14pt}
\begin{enumerate}[label=\textbf{\arabic*.}, leftmargin=*, itemsep=4pt, series=quiz]
\item What is the precession period of a Foucault pendulum placed at the North Pole? State your answer in sidereal days.



\vspace{2.0in}


\item If we double the length of a Foucault pendulum, what happens to its swing period? What happens to its precession period?
\vspace{2.0in}

\item The linearized equations of motion of the Foucault pendulum have the form $\mathbf{M}\,\ddot{\mathbf{q}} + \mathbf{C}\,\dot{\mathbf{q}} + \mathbf{K}\,\mathbf{q} = \mathbf{0}$,
where $\mathbf{C}$ is skew-symmetric ($\mathbf{C}^T = -\mathbf{C}$). What is the work done by the Coriolis force $-\mathbf{C}\,\dot{\mathbf{q}}$? Is the energy of the Foucault pendulum conserved?
\vspace*{\fill}

\end{enumerate}

\end{document}
